% texte.tex — f173r, vv. 876–919 (conversion initiale des travaux antérieurs)
\vers[n=877]{ΕΥ.}{Euripide}{ποίαν δὲ χώραν εἰσεκέλσαμεν σκάφει;}{Et quel est ce pays où ma nef a pris terre ?}
\vers[n=878]{ΜΝ.}{Mnésiloque}{Αἴγυπτον.}{L'Égypte.}
\vers[suite]{ΕΥ.}{Euripide}{ὦ δύστηνος οἷ πεπλώκαμεν.}{« Vers quel lieu, malheureux, nous avons navigué ! »}
\vers[n=879]{ΜΝ.}{Kritylla}{πείθει τι τῷ κακῷ κακῶς ἀπολουμένῳ}{T'en fies-tu à ce coquin — la peste soit de lui —}
\vers[n=880]{}{}{ληροῦντι λῆρον; Θεσμοφόριον τουτογί.}{qui te raconte des bourdes ? Ceci, c'est le Thesmophorion.}
\vers[n=881]{ΕΥ.}{Euripide}{αὐτὸς δὲ Πρωτεὺς ἔνδον ἔστ᾽ ἢ 'ξώπιος;}{Et lui, Protée, est-il céans, ou s'il est hors ?}
\vers[n=882]{ΚΡ.}{Kritylla}{οὐκ ἔσθ᾽ ὅπως οὐ ναυτιᾷς ἔτ᾽ ὦ ξένε,}{Il faut que tu aies encore mal au c\oe ur, étranger.}
\vers[n=883]{}{}{ὅστις γ᾽ ἀκούσας ὅτι τέθνηκε Πρωτέας}{On te dit qu'il est mort, Protéas,}
\vers[n=884]{}{}{ἔπειτ᾽ ἐρωτᾷς `ἔνδον ἔστ᾽ ἢ 'ξώπιος;'}{et tu demandes s'il est céans ou dehors.}
\vers[n=885]{ΕΥ.}{Euripide}{αἰαῖ τέθνηκε. ποῦ δ᾽ ἐτυμβεύθη τάφῳ;}{Hélas! il a vécu. Où donc est le sépulcre où il fut inhumé ?}
\vers[n=886]{ΜΝ.}{Mnésiloque}{τόδ᾽ ἐστὶν αὐτοῦ σῆμ᾽, ἐφ᾽ ᾧ καθήμεθα.}{C'est ici son tombeau, tu nous y vois assise.}
\vers[n=887]{ΚΡ.}{Kritylla}{κακῶς τ᾽ ἄρ᾽ ἐξόλοιο κἀξολεῖ γέ τοι,}{Puisses-tu périr misérablement ! Et tu périras,}
\vers[n=888]{}{}{ὅστις γε τολμᾷς σῆμα τὸν βωμὸν καλεῖν.}{toi qui oses appeler l'autel une tombe.}
\vers[n=889]{ΕΥ.}{Euripide}{τί δὴ σὺ θάσσεις τάσδε τυμβήρεις ἕδρας}{D'où vient qu'assise sur ce siège sépulcral,}
\vers[n=890]{}{}{φάρει καλυπτὸς ὦ ξένη;}{tu dérobes tes traits sous un voile, étrangère ?}
\vers[suite]{ΜΝ.}{Mnésiloque}{βιάζομαι}{On prétend sans mon v\oe u}
\vers[n=891]{}{}{γάμοισι Πρωτέως παιδὶ συμμεῖξαι λέχος.}{que le fils de Protée me reçoive en son lit.}
\vers[n=892]{ΚΡ.}{Kritylla}{τί ὦ κακόδαιμον ἐξαπατᾷς αὖ τὸν ξένον;}{Pourquoi, misérable, tromper encore cet étranger ?}
\vers[n=893]{}{}{οὗτος πανουργῶν δεῦρ᾽ ἀνῆλθεν ὦ ξένε}{Il est entré en fraude, ô étranger,}
\vers[n=894]{}{}{ὡς τὰς γυναῖκας ἐπὶ κλοπῇ τοῦ χρυσίου.}{auprès des femmes afin de toucher subrepticement à leur bijou.}
\vers[n=895]{ΜΝ.}{Mnésiloque}{κόνυζα τοὐμὸν σῶμα βάλλουσα ψόγῳ.}{Je suis la conyze qu'on accable d'un blâme.}
\vers[n=896]{ΕΥ.}{Euripide}{ξένη τίς ἡ γραῦς ἡ κακορροθοῦσά σε;}{Étrangère, dis-moi, qui donc est cette vieille, qui t'incrimine ainsi ?}
\vers[n=897]{ΜΝ.}{Mnésiloque}{αὕτη Θεονόη Πρωτέως.}{Elle ? C'est Théonoé, la fille de Protée.}
\vers[suite]{ΚΡ.}{Kritylla}{μὰ τὼ θεὼ}{Non, par les déesses,}
\vers[n=898]{}{}{εἰ μὴ Κρίτυλλά γ᾽ Ἀντιθέου Γαργηττόθεν·}{je ne suis autre que Critylla, fille d'Antithéos, du dème de Gargettos.}
\vers[n=899]{}{}{σὺ δ᾽ εἶ πανοῦργος.}{Toi, tu es un fripon.}
\vers[suite]{ΜΝ.}{Mnésiloque}{ὁπόσα τοι βούλει λέγε.}{Dis ce que tu voudras.}
\vers[n=900]{}{}{οὐ γὰρ γαμοῦμαι σῷ κασιγνήτῳ ποτὲ}{En épousant ton frère, je ne trahirai point}
\vers[n=901]{}{}{προδοῦσα Μενέλεων τὸν ἐμὸν ἐν Τροίᾳ πόσιν.}{Ménélas, mon époux, qui reste à Troie.}
\vers[n=902]{ΕΥ.}{Euripide}{γύναι τί εἶπας; στρέψον ἀνταυγεῖς κόρας.}{Qu'as-tu dit ? Femme, tourne vers moi ta brillante prunelle.}
\vers[n=903]{ΜΝ.}{Mnésiloque}{αἰσχύνομαί σε τὰς γνάθους ὑβρισμένη.}{J'ai honte, après l'outrage imprimé sur mes joues.}
\vers[n=904]{ΕΥ.}{Euripide}{τουτὶ τί ἔστιν; ἀφασία τίς τοί μ᾽ ἔχει.}{Qu'est-ce là ? Une stupeur me prend.}
\vers[n=905]{}{}{ὦ θεοὶ τίν᾽ ὄψιν εἰσορῶ; τίς εἶ γύναι;}{Ô dieux, que vois-je ? Femme, qui donc es-tu?}
\vers[n=906]{ΜΝ.}{Mnésiloque}{σὺ δ᾽ εἶ τίς; αὑτὸς γὰρ σὲ κἄμ᾽ ἔχει λόγος.}{Qui donc es-tu toi-même ? « C'est une question qui m'appartient aussi ».}
\vers[n=907]{ΕΥ.}{Euripide}{Ἑλληνὶς εἶ τις ἢ 'πιχωρία γυνή;}{« Es-tu Hellène, ou bien es-tu de ce pays ? »}
\vers[n=908]{ΜΝ.}{Mnésiloque}{Ἑλληνίς. ἀλλὰ καὶ τὸ σὸν θέλω μαθεῖν.}{« Hellène. Et à mon tour je veux savoir ta race. »}
\vers[n=909]{ΕΥ.}{Euripide}{Ἑλένῃ σ᾽ ὁμοίαν δὴ μάλιστ᾽ εἶδον γύναι.}{« Ô femme, je te vois tout pareille à Hélène. »}
\vers[n=910]{ΜΝ.}{Mnésiloque}{ἐγὼ δὲ Μενελέῳ σ᾽ ὅσα γ᾽ ἐκ τῶν λαιφίων.}{« Toi-même à Ménélas, si j'en juge par ces hardes.»}
\vers[n=911]{ΕΥ.}{Euripide}{ἔγνως ἄρ᾽ ὀρθῶς ἄνδρα δυστυχέστατον.}{« Oui, tu l'as reconnu, ce grand infortuné. »}
\vers[n=912]{ΜΝ.}{Mnésiloque}{ὦ χρόνιος ἐλθὼν σῆς δάμαρτος ἐσχάρας}{« Ô toi, le tard venu dans les bas de ta femme »,}
\vers[n=914, indent=2.5em, findent=1.5em, metre={do.}]{}{}{λαβέ με λαβέ με πόσι, περίβαλε δὲ χέρας.}{prends, prends-moi, mon époux. Presse-moi dans tes bras.}
\vers[n=915, indent=2.5em, findent=1.5em]{}{}{φέρε σὲ κύσω. ἄπαγέ μ᾽ ἄπαγ᾽ ἄπαγέ με}{Voyons, que je te baise. Emmène-moi, emmène, emmène-moi, et au plus vite.}
\vers[n=916]{}{}{λαβὼν ταχὺ πάνυ.}{(Il descend de l'autel, les femmes l'entourent.)}
\vers[suite]{ΚΡ.}{Kritylla}{κλαύσετ᾽ ἄρα νὴ τὼ θεὼ}{Ah! il pleurera, par les deux déesses,}
\vers[n=917]{}{}{ὅστις σ᾽ ἀπάξει τυπτόμενος τῇ λαμπάδι.}{celui qui t'emmènera. Je le frapperai de ma torche.}
\vers[n=918]{ΕΥ.}{Euripide}{σὺ τὴν ἐμὴν γυναῖκα κωλύεις ἐμέ,}{Quoi ? mon épouse à moi, fille de Tyndare,}
\vers[n=919]{}{}{τὴν Τυνδάρειον παῖδ᾽, ἐπὶ Σπάρτην ἄγειν;}{tu voudrais m'empêcher de la conduire à Sparte ?}
